% !TeX root = ../main.tex
\chapter{Wstęp}

\section{Uzasadnienie podjęcia tematu}

Trening klasyfikatora i detektora zwykle zapisuje się w PyTorch \cite{paszke2019pytorch}. W kodzie programu widać tensor oraz polecenie liczenia gradientu. Alokacja pamięci GPU, wybór algorytmu splotu i moment, w którym procesor czeka na urządzenie, zostają w runtime biblioteki. LibTorch przenosi ten sam sposób pracy do C++. Dla aplikacji, która tylko uczy model, taki podział jest wygodny. Dla pracy, która ma opisać koszt jednego kroku, jest niewystarczający: z samego wywołania propagacji wstecznej nie wynika, który bufor żyje przez całą epokę, a który powstaje na jeden operator.

W odpowiedzi na to powstał silnik DeepLearnLib, pisany w C++17 i uruchamiany na stosie CUDA. Tensor urządzenia dostaje pamięć z \api{cudaMalloc} i zwalnia ją razem z ostatnim właścicielem. Splot, redukcja maksymalizująca i normalizacja wsadowa są wywołaniami interfejsu C biblioteki cuDNN \cite{chetlur2014cudnn}. Warstwa gęsta oraz inne iloczyny macierzy idą przez \api{cublasGemmEx}. Leaky ReLU, dropout, aktualizacja wag w miejscu oraz strata YOLOv1 są jądrami przygotowanymi w tym repozytorium. cuDNN nie dostarcza straty detekcyjnej o układzie siatki YOLOv1, więc ta część musi być lokalna. Każda warstwa ma \api{forward}, \api{backward} i \api{step}. Krok uczenia idzie listą warstw od wejścia do głowicy, a potem od głowicy do wejścia. \api{Network::forward} w \api{src/Network.cpp} jest pętlą \api{for} po \api{layers\_}. Po każdym \api{forward} woła \api{view}. Pytanie pracy jest praktyczne: czy jawny właściciel bufora, jawne wywołanie cuDNN i jawna pętla wstecz wystarczą, żeby na trzech zadaniach uzyskać czas epoki oraz metrykę, które da się zestawić z LibTorch. Wstęp tego nie rozstrzyga. Liczby stoją w rozdziale~4.

W \api{src/Tensor.cpp} funkcja \api{Tensor::ensure} woła konstruktor tylko wtedy, gdy \api{std::optional} jest pusty albo \api{get\_shape}, \api{get\_device} albo \api{get\_dtype} różni się od argumentu. W drugiej gałęzi zwraca \api{*slot}. Konstruktor GPU woła \api{cudaMalloc}, więc powtórzony kształt tej alokacji nie wykonuje. Przestrzeń robocza \api{Conv2d} jest osobnym \api{optional} i przechodzi przez to samo wywołanie po pierwszym doborze algorytmu. \api{cudaStreamSynchronize} w tym pliku stoi w \api{to\_host} za \api{cudaMemcpyAsync} z urządzenia na host. Program woła \api{to\_host} przy skalarze zapisywanym do CSV, w \api{Network::save} i przy mAP. W pliku wyników kolumna zajętości powinna być wtedy płaska. Wzrost między epokami oznacza, że pętla jednak alokuje. Czas epoki pokazuje wtedy koszt liczenia i koszt dostarczenia wsadu. LibTorch pozostaje w repozytorium jako drugie binarium, o nazwie z przyrostkiem \api{_torch}. Oba programy czytają ten sam blok parametrów: rozmiar wsadu, liczbę epok i harmonogram współczynnika uczenia. Rdzeń \api{DeepLearnLib} nie linkuje LibTorch, więc pomiar własnego stosu nie zależy od tej biblioteki.

\section{Cel pracy}

Celem jest jeden framework do trzech zadań, złożony z tych samych warstw. Klasyfikacja obrazu używa sieci \api{SimpleCNN} i straty \api{CrossEntropyLoss}. Zbiory to MNIST oraz CIFAR-10 \cite{lecun1998mnist,krizhevsky2009cifar}. Wejście jest tensorem NCHW. CIFAR-10 jest czytany jako obrazy kolorowe, z bokiem 32, o ile plik JSON nie poda innego. MNIST pochodzi ze spakowanego pliku. Bok i liczba kanałów biorą się z jego nagłówka. Dodatkową metryką jest dokładność. Dane tabelaryczne obsługuje perceptron o dwóch warstwach gęstych, również ze stratą entropii krzyżowej. Konfiguracja obejmuje zbiór Iris, zbiór Wisconsin i mały zestaw demonstracyjny. Detekcja używa modelu YOLOv1 \cite{redmon2016yolo} oraz straty \api{YOLOLoss}. Obraz wejściowy ma bok 448. Programy obejmują PASCAL VOC z 20 klasami \cite{everingham2010voc}, zbiór BCCD z 3 klasami i zbiór syntetyczny, także z 3 klasami. Liczba klas zmienia szerokość ostatniej warstwy i szerokość straty. Metryką detekcji jest średnia precyzja przy progu IoU równym 0,5.

Wspólny jest kontrakt warstwy, a nie wspólna topologia wkompilowana w rdzeń. Skompilowana biblioteka \api{DeepLearnLib} zawiera tensor, warstwy, straty i loadery. Model YOLOv1 oraz \api{SimpleCNN} leżą obok rdzenia i trafiają do biblioteki \api{DeepLearnModels}. Perceptron tabelaryczny powstaje w programie treningowym przez złożenie dwóch warstw gęstych i Leaky ReLU. Nowy eksperyment zmienia ułożenie bloków albo plik konfiguracyjny. Alokacja, deskryptory cuDNN i krok optymalizatora zostają po stronie rdzenia. Współczynnik uczenia, momentum, rozkład wag i próg obcięcia gradientu są polami warstwy. Program treningowy przepisuje je na starcie epoki, więc splot i warstwa gęsta dostają ten sam krok. W katalogu \api{benchmarks} nie ma wywołania \api{Network::fit}. Pętlę epoki, harmonogram i zapis CSV mają pliki \api{train\_*\_custom.cpp}. Klasa \api{Network} udostępnia przebieg w przód, obcięcie gradientu oraz zapis parametrów. Nie ona prowadzi eksperymenty opisane w rozdziale~4.

Osobnym celem jest pomiar. Dla zbioru buduje się binarium własne. Gdy konfiguracja CMake znajduje LibTorch, buduje się też binarium referencyjne. Oba czytają blok z \api{config/experiments.json}. Po epoce program dopisuje wiersz CSV. Każda pętla detekcji zapisuje stratę uczenia, stratę testu, czas, zajętość pamięci urządzenia i mAP. Klasyfikacja obrazu zapisuje obie straty i obie dokładności. Iris i Wisconsin zapisują stratę uczenia i dokładność uczenia, bez części testowej. Skopiowane pliki powstałe wcześniejszym nagłówkiem nie mają części tych kolumn. Czas w tym pliku jest czasem zegarowym całej epoki, razem z ewaluacją. Czas jednego operatora mierzą mikrobenchmarki: jądra elementowe, iloczyn o kształcie głowicy, splot, normalizację wsadową oraz blok \api{FusedCBR2d}. Krok na danych VOC mierzą osobno \api{bench\_voc\_custom} i \api{bench\_voc\_torch}. Osobne programy inferencji wczytują zapisane parametry i liczą sam przebieg w przód. W detekcji ramki powstają na hoście, po progu ufności i po tłumieniu niemaksymalnym. W klasyfikacji decyzją jest indeks największego logitu.

\section{Teza}

Przyjęto tezę, że rdzeń C++ oparty na CUDA, cuBLAS i cuDNN uczy trzy rodziny modeli bez linkowania LibTorch: klasyfikator obrazowy \api{SimpleCNN}, dwuwarstwowy perceptron na danych tabelarycznych oraz detektor YOLOv1. Sprawdzenie wykonują pary programów na wspólnym protokole. Protokół obejmuje jedną kartę, rozmiar wsadu, liczbę epok i harmonogram współczynnika uczenia. Udziały podziału detekcji są wspólne. Lista plików nie jest, bo \api{split\_dataset} tasuje ją przy każdym uruchomieniu. Miarami są czas epoki mierzony zegarem ściennym, zajętość pamięci urządzenia oraz metryka zadania, czyli dokładność albo średnia precyzja przy progu IoU równym 0,5. Rozdział~4 zestawia te wielkości dla obu stosów. Na tej podstawie teza zostaje utrzymana albo ograniczona do tych zadań, na których rzędy czasu i metryki są zgodne.

\section{Zakres pracy i świadome ograniczenia}

Projekt składa się kompilatorem C++17, z zestawem CUDA w wersji co najmniej 12, razem z cuBLAS i cuDNN. Rdzeń obejmuje tensor gęsty w FP32. Warstwy to splot \api{Conv2d}, normalizacja wsadowa \api{BatchNorm2d}, blok \api{FusedCBR2d}, redukcja maksymalizująca \api{MaxPool2d}, warstwa gęsta \api{FullyConnected}, Leaky ReLU, dropout, spłaszczenie i softmax. Straty to \api{CrossEntropyLoss}, \api{MSELoss} i \api{YOLOLoss}. Na kartę trafiają cztery rodzaje danych: obrazy z adnotacją w układzie VOC (\api{CustomDataLoader}), obrazy ułożone w katalogach klas (\api{ClassificationLoader}), obrazy ze spakowanego pliku (\api{PackedImageLoader}) oraz wiersze CSV (\api{CSVLoader}). Klasa \api{Network} trzyma warstwy w kolejności przebiegu w przód, obcina gradient i zapisuje parametry. Kontekst cuBLAS i kontekst cuDNN są jedne na proces. \api{check\_cuda}, \api{check\_cublas} i \api{check\_cudnn} przy statusie innym niż sukces składają napis z \api{__FILE__} i \api{__LINE__} i rzucają \api{std::runtime\_error}. Makra \api{CHECK\_CUDA}, \api{CHECK\_CUBLAS} i \api{CHECK\_CUDNN} przekazują te argumenty.

Część aplikacyjna obejmuje modele, trening i inferencję. Detekcja ma pełny przebieg VOC, skrócony przebieg trzech epok, próbę przeuczenia na niewielkim wycinku VOC, trening BCCD oraz trening zbioru syntetycznego. Każdy z tych programów woła mAP. Przeuczenie liczy stratę testową i mAP na tym samym wycinku. Klasyfikacja obejmuje MNIST, CIFAR-10 oraz trzy zestawy tabelaryczne. Ramki do mAP powstają po kopiowaniu predykcji na host, po progu ufności i po tłumieniu niemaksymalnym. Sama średnia precyzja jest liczona na procesorze, w wariancie 11-punktowym, przy progu IoU równym 0,5. Oba progi biorą się z bloku JSON danego zbioru. Plik \api{config/sanity.json} ma ten sam schemat pól, lecz krótkie liczby epok. Służy do sprawdzenia, że program po kompilacji wykonuje krok na GPU. Czasy z tego pliku nie wchodzą do tabel pracy. Ciągła integracja kompiluje binaria własne na maszynie bez karty NVIDIA i sprawdza samo linkowanie. Wiersze czasu oraz pamięci pochodzą z lokalnego uruchomienia na karcie NVIDIA GeForce RTX 5080. Skrypt zbierający wyniki przepisuje ostatni wiersz istniejącego CSV i nie dopisuje pomiaru, którego plik nie zawiera.

Zakres jest ograniczony świadomie. Propagacja wsteczna jest zwykłą metodą warstwy. W \api{src/Network.cpp} nie ma typu grafu. \api{Network::forward} jest pętlą po \api{layers\_}. Uczenie działa na jednym urządzeniu CUDA. W \api{CMakeLists.txt} brak \api{CUDNN\_LIBRARY} kończy konfigurację wywołaniem \api{message(FATAL\_ERROR)}. W \api{src/Tensor.cpp} gałąź bez \api{DEEPLEARNLIB\_ENABLE\_CUDA} rzuca wyjątek z \api{matmul\_into} i z operatorów elementowych. Detektor w pracy to YOLOv1 w postaci zaimplementowanej w modelu aplikacyjnym, ze stratą na siatce używanej przez \api{YOLOLoss}. Późniejszych wersji YOLO nie składano. W zakresie są testy jednostkowe warstw oraz test zgodności numerycznej z operatorami LibTorch. Sam czas epoki nie rozstrzyga, czy oba stosy wykonały ten sam krok. LibTorch jest linkowany przez binaria referencyjne oraz przez te testy. Nie jest zależnością rdzenia. Przechowywanie na GPU jest FP32.

\section{Układ pracy}

Rozdział~2 zbiera operacje, z których składa się silnik: aktualizację wag, warstwy, funkcje straty, kodowanie ramki YOLOv1, metryki oraz wykonanie na GPU. Wzór pojawia się tylko wtedy, gdy ta sama postać wraca później w kodzie, i ma odsyłacz do źródła.

Rozdział~3 opisuje tensor, kontrakt warstwy oraz realizacje na cuDNN i cuBLAS. Następnie omawia loadery, złożenie trzech modeli z tych samych bloków oraz pętlę epoki zapisaną w programach treningowych.

Rozdział~4 podaje protokół porównania ze stosem LibTorch i wyniki klasyfikacji, danych tabelarycznych oraz treningu YOLOv1. Osobno rozdziela czas całej epoki od czasu pojedynczego operatora.

Rozdział~5 odpowiada na tezę i wymienia ograniczenia wynikające z architektury rdzenia. Załączniki zawierają skrót pól konfiguracji, układ pliku z wagami oraz tabele pomiarów zebrane w jednym miejscu.
